\subsection{Extension: enlarged training set (assumption A4)}\label{sec:enlarged}
Because the learning curves had not saturated and the assignment allows the number of samples to be increased, $30{,}000$ additional samples per case were drawn from the same distributions
(independent random stream) and used as extra \emph{training} data; the validation and test sets are the original ones, so all numbers are directly comparable. The comparison of the four model types
above is unchanged (it uses the specified $10{,}000$-row dataset); only the winning type, the MLP, was refit on $37{,}000$ training rows (1500 epochs; a first attempt with 500 epochs was under-trained and made the NMOS ON model worse).
\begin{table}[H]
\centering
{\small
\begin{tabular}{lrrrrr}
\toprule
case & train rows (enlarged) & median err (10k) & median err (enlarged) & 95th pct (10k) & 95th pct (enlarged) \\
\midrule
nmos\_on & 37000 & 4.05\% & 7.15\% & 22.72\% & 23.63\% \\
nmos\_off & 37000 & 1.62\% & 1.03\% & 20.18\% & 7.68\% \\
pmos\_on & 37000 & 1.91\% & 1.48\% & 11.32\% & 6.34\% \\
pmos\_off & 37000 & 4.23\% & 2.91\% & 22.66\% & 14.64\% \\
\bottomrule
\end{tabular}
}
\caption{Extension (assumption A4): test-set error of the MLP trained on the 7,000-row training split (headline) and refit on 37,000 rows (7,000 + 30,000 extra); same test set.}
\label{tab:enlarged}
\end{table}


For three of the four cases the enlarged training set lowers the test errors clearly (e.g.\ the 95th-percentile error of the NMOS OFF model falls from 20.2\,\% to 7.7\,\%, of the PMOS ON model from 11.3\,\% to 6.3\,\%, of the PMOS OFF model from 22.7\,\% to 14.6\,\%),
but the NMOS ON model does not improve: its median error is higher (7.2\,\% against 4.1\,\%) and it shows a systematic bias of several percent on the validation set.
Since the gate and circuit leakages are \emph{sums} of currents, a systematic bias matters more than the median error. The choice between the $10{,}000$-row and the enlarged model of each case was therefore made on the validation set only, with the score
$\tfrac12(\text{median}+\text{95th percentile})+2\,\overline{|\text{bias}|}$ (Table~\ref{tab:selection}); the bias term was added after we had seen that the enlarged NMOS ON model hurt the gate level, which is why this criterion is reported openly here.
It selects the $10{,}000$-row model for NMOS ON and the enlarged models for the other three cases (``enlarged-set selection'').
\begin{table}[H]
\centering
{\small
\begin{tabular}{llrrrr}
\toprule
case & candidate & val median err & val 95th pct & val mean |bias| & selection score \\
\midrule
nmos\_on & 10k & 4.30\% & 22.29\% & 0.79\% & 0.1487 \\
nmos\_on & enlarged & 7.19\% & 24.16\% & 6.11\% & 0.2789 \\
nmos\_off & 10k & 1.69\% & 14.82\% & 0.21\% & 0.08672 \\
nmos\_off & enlarged & 0.99\% & 7.17\% & 0.18\% & 0.04451 \\
pmos\_on & 10k & 2.04\% & 12.16\% & 0.36\% & 0.07812 \\
pmos\_on & enlarged & 1.54\% & 6.91\% & 0.09\% & 0.04405 \\
pmos\_off & 10k & 4.29\% & 24.75\% & 0.89\% & 0.1629 \\
pmos\_off & enlarged & 2.89\% & 14.76\% & 0.34\% & 0.09512 \\
\bottomrule
\end{tabular}
}
\caption{Validation-set selection between the 10,000-row model and the enlarged-training-set model (score = $\frac{1}{2}$(median + 95th pct) + 2 mean $|$bias$|$, lower is better).}
\label{tab:selection}
\end{table}



\paragraph{Effect at gate and circuit level.}
\begin{table}[H]
\centering
{\small
\begin{tabular}{llrrrr}
\toprule
models & level & bias & median err & 95th pct err & mean accuracy \\
\midrule
10,000-row models (headline) & gates (3,600 runs) & 1.51\% & 5.53\% & 19.56\% & 92.79\% \\
10,000-row models (headline) & parity checker (51,200 runs) & 0.96\% & 6.38\% & 19.12\% & 92.49\% \\
enlarged-set selection (extension) & gates (3,600 runs) & -0.52\% & 6.56\% & 22.62\% & 91.58\% \\
enlarged-set selection (extension) & parity checker (51,200 runs) & -0.16\% & 6.82\% & 19.63\% & 92.01\% \\
\bottomrule
\end{tabular}
}
\caption{Gate-level and circuit-level accuracy of the two model sets.}
\label{tab:sets}
\end{table}


The enlarged-set selection reduces the bias (from $+1.5\,\%$ to $-0.5\,\%$ at gate level) but does \emph{not} improve the gate-level accuracy: the median error is 6.6\,\% instead of 5.5\,\% and the 95th percentile 22.6\,\% instead of 19.6\,\%;
at circuit level the two sets are practically equal (6.8\,\% against 6.4\,\% median error). The better accuracy on the uniformly distributed validation samples therefore does not carry over to the operating points that occur in gates.
A possible reason, which was \emph{not} tested, is that the gate operating points lie on the edges of the sampled voltage range (inputs exactly at $0$ or $V_{DD}$, $V_s=0$, $V_d\approx V_{DD}$, ON devices at $V_{ds}\approx0$) whereas the validation error is averaged over interior points.
The headline results of this report are therefore those of the $10{,}000$-row models, i.e.\ the setting specified in the assignment; the enlarged-set models are reported as an extension.
